\chapter*{ТЕРМИНЫ И ОПРЕДЕЛЕНИЯ}
\addcontentsline{toc}{chapter}{ТЕРМИНЫ И ОПРЕДЕЛЕНИЯ}
\setlength{\parindent}{0pt}
% Перечень располагается столбцом, слева без абзацного отступа в алфавитном порядке,
% справа через тире — определение, без знаков препинания в конце строки (п. 4.9.2)
% Термины пишутся без жирного шрифта
\setlength{\LTpre}{0pt}
\setlength{\LTpost}{0pt}
\begin{longtable}{@{}P{4.8cm}@{\hspace{0.5em}—\hspace{0.5em}}P{10.4cm}@{}}
Source provenance & фиксация происхождения результата с указанием исходного документа, фрагмента источника, шага агентного конвейера и блока сгенерированного курса\\[0.5em]
Агентный конвейер & управляемая последовательность узлов ИИ-методиста, где каждый узел выполняет отдельный этап от анализа источника до финальной верификации\\[0.5em]
Знак качества & пороговый статус финальной верификации, рассчитанный по доле подтверждённых утверждений в таблице свидетельств\\[0.5em]
Таблица свидетельств & машиночитаемая таблица проверки фактических утверждений, содержащая утверждение, источник, статус подтверждения и комментарий проверки
\end{longtable}
\setlength{\parindent}{1.25cm}
